\documentclass[a4paper,12pt]{article}
\usepackage[utf8]{inputenc} % Kodowanie pliku
\usepackage[T1]{fontenc}    % Kodowanie czcionki
\usepackage[polish]{babel}  % Obsługa języka polskiego
\usepackage{graphicx}       % Obsługa obrazów
\usepackage{mathptmx}       % Times New Roman (podobna)
\date{ }
\usepackage{blindtext}
\usepackage{titlesec}
\usepackage{tocbibind}

\begin{document}
	
\begin{titlepage}
	\begin{center}
		% Logo
		\includegraphics[width=0.35\textwidth]{LOGO.png}
		\vspace{1cm}
		
		{\fontsize{20}{22}\selectfont \textbf{Politechnika Śląska}\\ \vspace{0.5cm} Wydział Automatyki, Elektroniki i Informatyki\vspace{1cm}}
		
		% Tytuł
		{\fontsize{18}{22}\selectfont \textbf{Projekt inżynierski}}
		
		
		\fontsize{14}{18}\selectfont "Gra planszowa 2.5D z uczeniem maszynowym komputerowych graczy"
		
		\vspace{1.5cm}
		
		% Autor
		{\fontsize{14}{18}\selectfont \textbf{Jagoda WIĘCEK}}
		\\Nr albumu: 300503
		
		\vfill
		
		{\fontsize{14}{18}\selectfont 
			\textbf{Kierunek:} Informatyka\\
			
			\textbf{Prowadzący pracę:} Dr inż. Ewa Lach
			\\Katedra Grafiki, Wizji Komputerowej i Systemów Cyfrowych
			\textbf{\\Gliwice 2024}
		
		}
		
		\vspace{0.8cm}
		
		
	\end{center}
\end{titlepage}
	

\tableofcontents
\newpage % Przejście na nową stronę

% Sekcje dokumentu
\section{Streszczenie} %TODO
TODO
\section{Wstęp} %DONE
Gry planszowe są formą rozrywki dla ludzi od wieków jak i formą intelektualnego wyzwania. We współczesnych czasach znaczna część życia podlega cyfryzacji, od komunikacji do płatności bankowych. Z tego powodu również wszelkie gry planszowe są przenoszone do cyfrowego świata. Takie rozwiązanie pozwala na rozpoczęcie rozgrywki ulubionej gry planszowej w świecie online wraz z graczami z całego świata bez żadnych ograniczeń. Są to jednak głównie gry online, zakładają one grę z żywym przeciwnikiem, nie rozgrywkę z maszyną. Ten projekt zakłada możliwość napisania gry, w której gracz będzie mieć do czynienia z graczem komputerowym wraz z zaimplementowanym algorytmem uczenia maszynowego.
\subsection{Cel pracy} 
Celem pracy jest zaimplementowanie gracza komputerowego z algorytmami nauczania maszynowego w unity. Zamierzonym efektem tego projektu ma być gracz komputerowy, który będzie w stanie wykonywać ruchy na planszy jak żywy gracz. 
\subsection{Zakres pracy}
W projekcie należy uwzględnić: 
\\ - wybraną grę planszową oraz jej zasady
\\ - sposób poruszania się dla gracza
\\ - ograniczenia i możliwe ruchy dla modelu
\\ - kary i nagrody po wykonaniu konkretnych ruchów


\subsection{Charakterystyka rozdziałów}
\textbf{Wstęp} - przedstawienie celu oraz zakres pracy, określenie grup docelowych
\\ \textbf{Analiza tematu} - zawiera ogólną analizę wszystkich możliwych rozwiązań tego projektu.
\\ \textbf{Wymagania i narzędzia} - przedstawienie wymagań, jakie projekt narzucił oraz wykorzystane narzędzia do jego wykonania
\\ \textbf{Specyfikacja zewnętrzna} - określenie wymagań sprzętowych, scenariusz korzystania z aplikacji oraz instalacja potrzebnych narzędzi. 
\\ \textbf{Specyfikacja wewnętrzna} - przedstawienie wykorzystanych bibliotek w projekcie, opis najważniejszych klas oraz diagramy.
\\ \textbf{Weryfikacja i walidacja} - przedstawienie możliwych modyfikacji dla modelu, które mogą wpływać na tok nauczania, testowanie aplikacji pod kątem nieprzewidzianych błędów. 
\\ \textbf{Podsumowanie i wnioski} - obserwacje toku nauki modelu i wnioski względem możliwości rozwoju.

\section{Analiza tematu}%mniej więcej zrobione
Tematem pracy jest cyfrowa gra planszowa z wykorzystaniem algorytmów sztucznej inteligencji. Kluczowym problemem jest zaimplementowanie inteligentnego gracza, który poprzez system prób i błędów, nauczy się grać. Wybranie gry planszowej ma duże znaczenie, im bardziej skomplikowane zasady gry, tym trudniej będzie modelowi nauczyć się zasad.
\\Do realizacji tego projektu została wybrana gra planszowa o nazwie Rummikub.
\subsection{Sformułowanie problemu}
Gry planszowe w wersji cyfrowej są kodowane w celu gry online z innymi osobami za pośrednictwem internetu. Jednak zaprogramowanie gry z komputerowo uczącym się graczem wymaga innej taktyki, model nie może się uczyć na graniu z żywymi graczami, gdyż zajęło by mu to za długo. Algorytmy AI potrzebują dużej ilości powtórzeń aby nauczyć się zasad gry. W projekcie zostanie uwzględniony jeden żywy gracz, gracze komputerowi z predefiniowanymi ruchami oraz komputerowy gracz z algorytmem nauczania maszynowego.
\subsection{Znane rozwiązania}
Na rynku jest wiele gier planszowych, najpopularniejsze z nich doczekały się swojej cyfrowej wersji. Są nimi między innymi:\\
	- Monopoli\\
	- Cantan\\
 	- Rummikub
 - Wydana gra online Rummikub jest nastawiona na grę z użyciem internetu do komunikacji innych graczy z całego świata. Gra jest zaimplementowana w grafice 2D bez algorytmów sztucznej inteligencji. 

\section{Wymagania i narzędzia}%TODO
W tej sekcji przedstawiono wymagania funkcjonalne i niefunkcjonalne projektu oraz narzędzia technologiczne i programistyczne, które zostały wykorzystane do jego realizacji. Opisane zostaną także kluczowe zasoby niezbędne do wdrożenia rozwiązania, takie jak środowiska programistyczne, biblioteki, frameworki oraz sprzęt.
\subsection{Wymagania funkcjonalne}
\textbf{Interakcja z mapą :}
\begin{itemize}
\item Gra musi pozwolić graczowi i agentowi na wybór płytki, którą chcą postawić bądź wejść z nią w interakcje.
\item Gracz oraz agent powinni być w stanie postawić nową płytkę na mapie, przesunąć istniejącą bądź usunąć ją z mapy, jeśli jest to możliwe.
\item Agent oraz gracz muszą układać sekwencje płytek jak najszybciej.
\item Gra musi ograniczać graczy, aby nie mogli tworzyć sekwencji niedozwolonych w zasadach gry.
\end{itemize}
\textbf{Ruchy gracza :}
\begin{itemize}
	\item Użytkownik musi mieć możliwość widzieć swoje płytki i wykładać je na mapę oraz tworzyć 
	\item W trakcie tury gra na bieżąco weryfikuje czy daną płytkę można postawić, pod koniec tury sprawdza czy każda wyłożona sekwencja posiada co najmniej 3 płytki w jednym ciągu.
	\item Gracz powinien móc zakończyć turę, cofnąć wszystkie akcje oraz dobrać nową płytkę.
	\item Użytkownik powinien móc posortować swoje płytki seryjnie aniżeli ręcznie.
\end{itemize}
\subsection{Wymagania niefunkcjonalne}
\subsection{Opis narzędzi}
W tej sekcji przedstawiono narzędzia, które zostały wykorzystane do zrealizowania projektu.
\subsubsection{unity}
\subsubsection{C\#}
\subsubsection{ML-Agents}


\section{Specyfikacja zewnętrzna}%TODO
TODO
\subsection{Wymagania sprzętowe i programowe}
\subsection{Instalacja}
\subsubsection{unity}
\subsubsection{biblioteka ml-agents}
\subsection{Aktywacja}
\subsection{Sposób obsługi}
\subsection{Ekran aplikacji}
\subsubsection{Zasady gry}
\subsubsection{Opis ekranu gry}

\section{Specyfikacja wewnętrzna}%TODO
TODO
\subsection{Biblioteka ml-agents}
\subsection{Klasa PlayerIA}
\subsection{Klasa GameController}
\subsection{Klasa Tile}


\section{Weryfikacja i walidacja}%TODO
TODO
\subsection{Testowanie}
\subsection{Wykryte błędy}

\section{Podsumowanie i wnioski}%TODO
TODO
\subsection{Napotkane problemy}%TODO
\subsection{Możliwy rozwój}%TODO

\section{Bibliografia}%TODO
TODO
\section{Spis skrótów i symboli}%TODO
\section{Spis rysunków}%TODO

\end{document}%closing the document